\documentclass[../../../include/open-logic-section]{subfiles}

\begin{document}

\olfileid{sth}{story}{predicative}	

\olsection{محمولی و غیرمحمولی}

مجموعهٔ راسل، یعنی~$R$، به‌وسیلهٔ
$\Setabs{x}{x \notin x}$ تعریف شد. اگر این تعریف را به‌تفصیل بیان کنیم،
$R$ مجموعه‌ای خواهد بود که همه و تنها مجموعه‌هایی را دربر می‌گیرد که عضو
خود نیستند. پس در تعریف~$R$، بر دامنه‌ای سور می‌بندیم که (اگر~$R$ وجود
می‌داشت) خود~$R$ را نیز دربر می‌گرفت.

این تعریفی \emph{غیرمحمولی} است. به‌طور کلی‌تر، می‌توانیم بگوییم یک تعریف
غیرمحمولی است اگر و تنها اگر بر دامنه‌ای سور ببندد که شیء در حال تعریف را
دربر می‌گیرد.
	
پس از ظهور پارادوکس‌ها، وایتهد، راسل، پوانکاره و وایل چنین تعریف‌های
غیرمحمولی‌ای را به‌سبب «دور باطل» رد کردند:
\begin{quote}
	تحلیل پارادوکس‌هایی که باید از آن‌ها پرهیز کرد نشان می‌دهد که همگی از
	نوعی دور باطل ناشی می‌شوند. دورهای باطل مورد بحث از این فرض پدید
	می‌آیند که گردایه‌ای از اشیا می‌تواند اعضایی داشته باشد که تنها به کمک
	آن گردایه به‌مثابه یک کل قابل تعریف‌اند[\ldots. \textparagraph]

	اصلی را که به ما امکان می‌دهد از تمامیت‌های نامجاز بپرهیزیم می‌توان
	چنین بیان کرد: «هر آنچه همهٔ یک گردایه را درگیر می‌کند، نباید یکی از
	اعضای آن گردایه باشد»؛ یا برعکس: «اگر گردایه‌ای معین، به فرض آنکه
	تمامیتی می‌داشت، اعضایی می‌داشت که تنها بر حسب آن تمامیت قابل تعریف
	بودند، آن گردایه هیچ تمامیتی ندارد.» این را «اصل دور باطل» خواهیم
	نامید، زیرا ما را قادر می‌سازد از دورهای باطلی که فرض تمامیت‌های نامجاز
	در پی دارد اجتناب کنیم.
	\citep[p.~37]{WhiteheadRussell1910}
\end{quote}
اگر در پذیرش \emph{اصل دور باطل} از آنان پیروی کنیم، ممکن است بکوشیم
طرحوارهٔ فاجعه‌بار تصریح ساده‌انگارانه (در~\olref[sth][story][rus]{sec})
را با چیزی مانند اصل زیر جایگزین کنیم:

\begin{defish}
\emph{تصریح محمولی.} برای هر فرمول~$\phi$ که تنها بر مجموعه‌ها سور می‌بندد، مجموعهٔ$^\prime$~$\Setabs{x}{\phi(x)}$ وجود دارد.
\end{defish}

تا هنگامی که مجموعه‌های$^\prime$ مجموعه نباشند، هیچ تناقضی پدید نخواهد آمد.

متأسفانه، تصریح محمولی چندان \emph{فراگیر} نیست. زیرا ما را با موجودات
جدیدی، یعنی مجموعه‌های$^\prime$، آشنا می‌کند. پس ناچار خواهیم بود
فرمول‌هایی را بررسی کنیم که بر مجموعه‌های$^\prime$ سور می‌بندند. اگر این
فرمول‌ها همواره مجموعه‌ای$^\prime$ به دست دهند، پارادوکس راسل دوباره رخ
خواهد داد؛ کافی است مجموعهٔ$^\prime$ همهٔ مجموعه‌های$^\prime$ خودناعضو را
در نظر بگیریم. بنابراین، اگر همین اندیشه را دنبال کنیم، باید بگوییم
فرمولی که بر مجموعه‌های$^\prime$ سور می‌بندد مجموعهٔ متناظر
$^{\prime\prime}$ را به دست می‌دهد. سپس به مجموعه‌های
$^{\prime\prime\prime}$، مجموعه‌های$^{\prime\prime\prime\prime}$ و جز آن
نیاز خواهیم داشت. برای جلوگیری از انبوه علامت‌های پرایم، آسان‌تر است آن‌ها
را مجموعه‌های$_0$، مجموعه‌های$_1$، مجموعه‌های$_2$، مجموعه‌های$_3$،
مجموعه‌های$_4$، \ldots{} بنامیم. از این راه به نظریهٔ (سادهٔ) انواع
می‌رسیم.

چند ایراد آشکار بر چنین نظریه‌ای وارد است (هرچند آشکار نیست که این ایرادها
\emph{قاطع} باشند). به‌اختصار: کار با نظریهٔ حاصل دشوار است؛ این نظریه در
فرض‌کردن اقسام متفاوت اشیا دست‌ودلباز است؛ و سرانجام روشن نیست که
تعریف‌های غیرمحمولی واقعاً \emph{تا این اندازه بد} باشند.
	
برای روشن‌شدن نکتهٔ اخیر، این سخن~\citeauthor{Ramsey1925} را در نظر
بگیرید:
\begin{quote}
	می‌توانیم از مردی به‌عنوان بلندقدترین فرد یک گروه یاد کنیم و بدین‌سان
	او را به کمک تمامیتی مشخص سازیم که خودش نیز عضو آن است، بی‌آنکه هیچ
	دور باطلی در کار باشد. \citep{Ramsey1925}
\end{quote}
نکتهٔ رمزی این است که «بلندقدترین مرد گروه» تعریفی \emph{غیرمحمولی}
است، اما آشکارا کاملاً بی‌اشکال است.

ممکن است پاسخ داده شود که در این مورد می‌توانیم بلندقدترین شخص را به
شیوه‌ای \emph{محمولی} مشخص کنیم؛ برای نمونه، شاید بتوانیم صرفاً به مرد
مورد نظر اشاره کنیم. در این صورت، ایراد به تعریف‌های غیرمحمولی آشکارا
باید به موجوداتی محدود شود که \emph{تنها} به‌نحوی غیرمحمولی قابل تشخیص‌اند.
اما حتی در آن صورت نیز لازم است بیشتر توضیح داده شود که چرا چنین
«غیرمحمولی‌بودن ذاتی‌ای» تا این اندازه بد است.\footnote{برای بحث بیشتر،
بنگرید به~\citet{Linnebo2010}.}

البته اگر بخواهیم تعریف‌هایمان چیزی مانند دستورالعمل \emph{ساختن} یک شیء
به ما بدهند، تعریف‌های غیرمحمولی بسیار مسئله‌سازند. زیرا با داشتن تعریفی
غیرمحمولی، واقعاً در دوری باطل گرفتار می‌شویم: برای ساختن شیئی که به‌نحوی
غیرمحمولی مشخص شده است، باید \emph{ابتدا} همهٔ اشیا (از جمله همان شیء
غیرمحمولی) را ساخته باشیم، زیرا آن شیء بر حسب همهٔ اشیا مشخص شده است؛ پس
باید آن شیء را پیش از آنکه خود آن را ساخته باشیم، بسازیم. اما باز هم این
تنها هنگامی ایرادی جدی به مجموعه‌هایی است که «به‌نحوی ذاتاً غیرمحمولی»
مشخص شده‌اند که مجموعه‌ها را چیزهایی بدانیم که خود \emph{می‌سازیم}. و
(احتمالاً) چنین نمی‌اندیشیم.

ازاین‌رو---خواه به سود ما باشد خواه به زیانمان---رویکردی که رایج شد در
پیش‌گرفتن موضعی سخت‌گیرانه دربارهٔ محمولی یا غیرمحمولی‌بودن نبود. در عوض،
این رویکرد از چیزی بهره می‌گیرد که اکنون رویکرد انباشتی-تکراری نامیده
می‌شود. این رویکرد سرانجام به ما امکان می‌دهد مجموعه‌هایمان را در
«مرحله»هایی لایه‌بندی کنیم---\emph{تا اندازه‌ای} همان‌گونه که رویکرد
محمولی موجودات را به مجموعه‌های$_0$، مجموعه‌های$_1$، مجموعه‌های$_2$،
\ldots{} لایه‌بندی می‌کند---اما هیچ تفاوتی در نوع آن‌ها فرض نخواهیم کرد.

\end{document}
